% =====================================================================
\chapter{METHODOLOGY}
\label{ch:method}
% =====================================================================

This chapter compares the candidate methods of Section~\ref{sec:candidates}
component by component and selects one for each
(Section~\ref{sec:selection}), then presents the design of the proposed
system: its structure, conceptual design and the justification of the
techniques used (Section~\ref{sec:design}).

\section{Approaches Comparison and Selection}
\label{sec:selection}

\subsection{Selection criteria}

The candidates are compared against six criteria derived from the design
requirements of Table~\ref{tab:requirements}:
\begin{description}[leftmargin=1.2cm,style=nextline,itemsep=2pt]
\item[C1] \emph{Truthfulness}: does the method keep reporting the true value
of time a dominant strategy (R1)?
\item[C2] \emph{Efficiency within the range}: is the selected allocation
optimal among the generated routes (R4)?
\item[C3] \emph{Bid independence}: is the method's output fixed before any
report is read (R5)?
\item[C4] \emph{Completeness or safety}: does a generation method find every
needed route, and does a reduction rule leave all outcomes unchanged (R7)?
\item[C5] \emph{Computational cost}: how much time and memory does it need at
the instance sizes of the study (R8)?
\item[C6] \emph{Verifiability}: can it be checked against an independent
oracle (R9)?
\end{description}
C1, C3 and C4 are knock-out criteria: a method that fails one of them
cannot be used, however fast it is.

\subsection{Mechanism}

Table~\ref{tab:cmp-mech} compares the candidate mechanisms. Posted prices
and pay-as-bid fail C1 or C2: under pay-as-bid drivers gain from inflating
their reports, and posted prices give up efficiency because the platform
cannot see costs. Approximate VCG-like mechanisms give up either exact
efficiency within the range or deterministic truthfulness. Exact VCG is the
only candidate that satisfies all knock-out criteria; its cost is one extra
exact solve per winner, which is acceptable if the solves are small, and
reducing their size is exactly what the pruning rules address.
\textbf{Selected: exact VCG with Clarke-pivot payments.} Posted prices and
pay-as-bid are kept as benchmarks for RQ3.

\begin{table}[t]
\centering
\caption{Comparison of candidate mechanisms}
\label{tab:cmp-mech}
\small
\begin{tabular}{P{2.6cm}P{2.5cm}P{2.1cm}P{2.6cm}P{3.6cm}}
\toprule
Mechanism & C1 truthful & C2 efficient in range & C5 cost & Main drawback\\
\midrule
Posted price & Accept/decline only & No & One solve & Price must be tuned; inefficient allocation\\
Pay-as-bid & No & On reports only & One solve & Drivers inflate reports; outcome depends on strategic reasoning\\
Approximate / greedy VCG & In expectation or approx. & No & Low & Loses exactness or deterministic truthfulness\\
\textbf{Exact VCG} & \textbf{Yes (DSIC)} & \textbf{Yes} & One solve per winner extra & Needs exact solves; payout premium (information rent)\\
\bottomrule
\end{tabular}
\end{table}

\subsection{Route generation}

Route generation was the subject of six preliminary verification rounds
(Tests~2--6 of the project), each of which checked a candidate against
brute-force enumeration on small instances ($n\le 6$, $B\le 4$) and measured
its cost. Table~\ref{tab:cmp-gen} summarises the evidence.

\begin{table}[t]
\centering
\caption{Comparison of candidate route-generation methods (preliminary verification rounds)}
\label{tab:cmp-gen}
\small
\begin{tabular}{P{3.4cm}P{1.7cm}P{4.6cm}P{4.0cm}}
\toprule
Method & C3 bid-indep. & C4 completeness (evidence) & C5 cost (evidence)\\
\midrule
Brute force & Yes & Complete by construction & Explosive; used only as test oracle\\
Level-wise insertion & Yes & Complete: 0 mismatches on 1,152 instances & Sequences per bundle reach the ceiling $(2k)!/2^k$ (2,520 at $k=4$) under wide windows\\
Insertion + Pareto dominance on sequences & Yes & \textbf{Fails}: 31.26\% of gigworker bundles lose feasible sequences & --\\
Insertion + lookahead-profile dominance & Yes & Sound at depth 1; \textbf{fails} at depth $\ge2$ (66.5--72.4\% violations) & Good compression at depth 1 only\\
Self-survival bound & Yes & Sound (0 violations on 62,728 checks) & Prunes 0\% where windows are wide\\
\textbf{Forward label-setting, exact key} & \textbf{Yes} & \textbf{0 violations on 616 instances; 0/451 and 0/493 in adversarial re-checks} & \textbf{5--10$\times$ faster than level-wise insertion for $n$ up to 15}\\
Column generation & \textbf{No} & -- & --\\
Detour-monotone branching & Yes & Complete & Applies to occasional drivers only\\
\bottomrule
\end{tabular}
\end{table}

The level-wise insertion method is complete but stores too many sequences.
Every attempt to compress it by comparing \emph{finished} sequences of the
same bundle through a few summary numbers failed, for a structural reason:
two sequences with the same summary can end at different positions, and a
later order may be reachable cheaply from one but not from the other (a
``hidden diagonal shortcut''). Forward label-setting avoids this flaw
because it compares only labels that share the exact current node, the set
of orders on board and the set of delivered orders; such labels have
exactly the same future options, and the dominance argument of
\citet{dumas1991pickup} applies. It passed every completeness check,
including hand-built two-dimensional counterexamples designed to defeat
it, and its compression improves as $n$ grows. Column generation is
excluded by C3, because its pricing step uses dual values that depend on the
reported costs. Detour-monotone branching is valuable but covers only one of
the two classes. \textbf{Selected: forward label-setting with exact-key
dominance for both classes (Algorithm~A), followed by a same-bundle Pareto
filter.}

\subsection{Pool reduction}

Table~\ref{tab:cmp-prune} compares the pool-reduction rules. Two rules that
look natural are unsafe. A per-driver hull across bundles deletes a route
whenever another route of the same driver is cheaper in both cost
components; in a synthetic audit it deleted 1,144 of 1,157 frontier routes,
and Example~\ref{ex:twin} shows that deleting such a route can change the
optimum. Merging labels under a coarser key (dropping the delivered set, or
keeping only its size) disagreed with brute-force enumeration in 46--47\% of
720 small instances (Section~\ref{sec:sd-algA} explains why). Instance-specific reduction by
sampling reports is not safe for unsampled reports, and any rule that reads
reports breaks C3. A non-local exact rule would be safe, but deciding it is
NP-hard (Remark~\ref{rem:scope}).

\begin{table}[t]
\centering
\caption{Comparison of candidate pool-reduction rules}
\label{tab:cmp-prune}
\small
\begin{tabular}{P{3.3cm}P{1.6cm}P{1.6cm}P{3.1cm}P{4.2cm}}
\toprule
Rule & C3 bid-indep. & C4 safe & Strength & Remark\\
\midrule
No reduction & Yes & Yes & None & Baseline\\
Same-bundle Pareto & Yes & Yes & Weak & Never compares different bundles\\
Per-driver hull across bundles & Yes & \textbf{No} & Strong & Deleted 1,144 of 1,157 frontier routes in an audit\\
Coarser label key & Yes & \textbf{No} & Strong & 46--47\% of 720 instances disagree with brute force\\
Report sampling & \textbf{No} & \textbf{No} & Very strong & Keeps only 0.4--1.5\% of the pool, but only for sampled reports\\
Non-local exact & Yes & Yes & Strongest & NP-hard to decide\\
\textbf{Local FD-completion frontier} & \textbf{Yes} & \textbf{Yes (proved)} & \textbf{Strong} & Post-processing; minimality planned\\
FD-completion label rule & Yes & Yes (in audit) & Reaches the same frontier & Planned; saves label extensions\\
\bottomrule
\end{tabular}
\end{table}

The local frontier is the only candidate that is bid-independent, provably
safe (Theorem~\ref{thm:frontier}) and still removes a large part of the pool.
\textbf{Selected: the local FD-completion frontier $\Pi^\star$ as a
post-processing step}, with the FD-completion label rule as a planned
accelerator that reaches the same frontier.

\subsection{Payment computation}

Naive re-solving is exact, simple and verifiable, and on the instance sizes
of this study its cost is small compared with route generation (at most
0.66 seconds median per instance at $n=20$). Component decomposition is
also exact, but only when the outside option is separable
(Section~\ref{sec:sd-decomp}); it is planned as a variant. Approximate counterfactuals fail C1
and C2. \textbf{Selected: naive exact re-solve as the reference method, with
component decomposition as an exact alternative under a separable outside
option.}

\subsection{Selected approach}

The selected approach is therefore an exact VCG mechanism on a
bid-independent route pool generated by forward label-setting, reduced to
the local FD-completion frontier, with winner determination and
counterfactual problems solved exactly as set-partitioning MILPs.

\section{Proposed Solution Approach}
\label{sec:design}

\subsection{System structure: inputs, process and outputs}

Figure~\ref{fig:pipeline} shows the structure of the proposed system, and
Table~\ref{tab:ipo} lists its inputs, processing stages and outputs.
Everything to the left of the report-collection step uses public data only;
this is what makes the allocation range independent of the reports.

\begin{figure}[t]
\centering
\resizebox{\textwidth}{!}{%
\begin{tikzpicture}[
  node distance=5mm and 7mm,
  every node/.style={font=\small},
  box/.style={rectangle, rounded corners, draw, align=center, minimum height=13mm, inner sep=2mm, fill=gray!6},
  arr/.style={-{Stealth[length=2mm]}, thick}]
\node[box] (data) {Public data\\$O$, $q$, $B$, $\Theta$,\\driver data};
\node[box, right=of data] (A) {Algorithm A\\route generation\\label dominance,\\Pareto per bundle};
\node[box, right=of A] (K) {Frontier $\Pi^\star$\\$R_i\mapsto\Kstar_i$\\certificate,\\pool hash};
\node[box, right=of K] (bids) {Collect reports\\$b\in\Theta^N$};
\node[box, right=of bids] (B) {Algorithm B\\winner\\determination};
\node[box, right=of B] (C) {Algorithm C\\Clarke-pivot\\payments};
\draw[arr] (data) -- (A); \draw[arr] (A) -- (K); \draw[arr] (K) -- (bids);
\draw[arr] (bids) -- (B); \draw[arr] (B) -- (C);
\begin{scope}[on background layer]
\node[draw, dashed, rounded corners, inner sep=3mm, fit=(data)(A)(K),
      label={[font=\small]above:pre-auction phase: no report is read}] {};
\node[draw, dashed, rounded corners, inner sep=3mm, fit=(B)(C),
      label={[font=\small]above:clearing phase}] {};
\end{scope}
\end{tikzpicture}}
\caption{Structure of the proposed procurement system. Route generation and
frontier pruning use public data only and are completed before reports are
collected; the pruned pools define the allocation range on which exact VCG
is run.}
\label{fig:pipeline}
\end{figure}

\begin{table}[t]
\centering
\caption{Inputs, process and outputs of the proposed system}
\label{tab:ipo}
\small
\begin{tabular}{P{2.1cm}P{12.3cm}}
\toprule
Element & Content\\
\midrule
Inputs & Orders (pickup and delivery locations, demand, release time, time windows); fixed-fleet prices $q_o$; bundle cap $B$; report domain $\Theta$; public driver data (start location, availability, capacity, distance coefficient $\kappa_i$; for occasional drivers also origin, destination, direct trip and detour budget); travel distance and time matrices; one report $b_i$ per driver.\\
\addlinespace
Process & (1) Algorithm~A generates each driver's route pool from public data. (2) The frontier rule $\Pi^\star$ deletes routes that cannot matter for any admissible report and certifies drivers with an empty frontier. (3) The pools are hashed and published; reports are collected. (4) Algorithm~B solves the winner-determination MILP exactly. (5) Algorithm~C solves one counterfactual MILP per winner and computes the payments.\\
\addlinespace
Outputs & Allocation (which driver serves which route, which orders go to the fixed fleet); payment to every driver; optimal value; pre-auction certificate (drivers with an empty frontier); pool hash; run statistics (pool sizes, times, solver status).\\
\bottomrule
\end{tabular}
\end{table}

\subsection{Conceptual design}

The conceptual design separates \emph{what the platform knows when}. The
platform holds three kinds of information (Table~\ref{tab:publicprivate}):
public order data and prices, public or committed driver data, and each
driver's private value of time. The system is organised into three phases
that respect this separation.

\begin{table}[t]
\centering
\caption{Public and private information in the proposed system}
\label{tab:publicprivate}
\small
\begin{tabular}{P{6.3cm}P{1.4cm}P{1.4cm}P{4.7cm}}
\toprule
Item & GW & OD & Status\\
\midrule
Start location, availability, capacity & Yes & Yes & Public\\
Distance coefficient $\kappa_i$ & Yes & Yes & Public or contractual\\
Personal destination, direct trip & -- & Yes & Public or committed before bidding\\
Detour budget $\tau_i$ & -- & Yes & Public or committed before bidding\\
Value of time $\theta_i$ & Yes & Yes & \textbf{Private}\\
Report $b_i$ & Yes & Yes & The only message sent\\
\bottomrule
\end{tabular}
\end{table}

\begin{enumerate}[leftmargin=*,itemsep=4pt]
\item \emph{Pre-auction phase.} From public data only, Algorithm~A builds
each driver's route pool and the frontier rule reduces it. The reduced pool
of each driver is described by a list of routes with their bundle, visiting
sequence and the two coefficients $(K,W)$ of the cost model. A driver whose
frontier is empty although she has a feasible route is certified as
dominated by the fixed fleet at every admissible report; if every driver
is certified, no auction needs to be opened. The pools are hashed, so that
drivers can later check that the allocation range was not changed after
bidding.
\item \emph{Auction phase.} Each remaining driver submits one number, her
reported value of time $b_i\in\Theta$.
\item \emph{Clearing phase.} The reported cost of every route is
$K+b_iW$. Algorithm~B selects routes and fixed-fleet assignments at minimum
total reported cost; Algorithm~C removes each winner in turn, re-solves and
pays the Clarke-pivot payment.
\end{enumerate}

The design reflects the concepts of Section~\ref{sec:litreview}: bid
independence (phase 1 never reads a report), exactness (phase 3 uses proven
optimal solutions), locality (the frontier of a driver uses only her own
pool and the public prices), the fixed fleet as a priced outside option
(it appears in the envelope that defines the frontier), and a single private
parameter (each route is summarised by two coefficients, so its cost is a
line in the report).

\subsection{Justification of the techniques}

\paragraph{Set-partitioning MILP and an exact solver.} Winner
determination is a set-partitioning problem with one convexity row per
driver. Exactness is a correctness requirement, not a quality target:
the truthfulness argument of Proposition~\ref{prop:dsic} needs proven
optimal solutions with report-independent tie-breaking. A commercial MILP
solver (CPLEX~12.10) with zero gap tolerance provides proven optimality; on
the pools of this study each solve takes a fraction of a second.

\paragraph{Label-setting for route generation.} Forward label-setting
builds routes in the order in which they are driven, so waiting, release
times and time windows are evaluated exactly at every step; its dominance
rule is safe because it compares only states with identical futures. It
handles the open routes of gigworkers and the destination-bound routes of
occasional drivers in one procedure.

\paragraph{The local frontier.} The frontier is the strongest reduction
that can be certified without reading reports and without reasoning about
other drivers, and it is exact for every admissible report. Because the
cost of each route is a line in the report, membership in the frontier is
decided by comparing finitely many lines, which is cheap.

\paragraph{Oracle-based verification and pre-registration.} Because every
claim of the mechanism rests on exactness, each component is checked against
an independent brute-force oracle on small instances before any main run,
and the theoretical results are additionally checked by a reference
implementation that shares no code with the production one, including
deliberate mutations that must be detected. All experimental parameters are
locked and hashed before the main runs, so that no parameter can be tuned
after results are seen.
